% ============================================================================
% ETHICAL, LEGAL, AND ENVIRONMENTAL CONSIDERATIONS
% ============================================================================
% Replace the existing section of the same name in the Discussion chapter with
% this block. It assumes that the bibliography contains the three entries given
% after the block.
% ============================================================================

\section{Ethical, legal, and environmental considerations}

The Index\_Insight pipeline processes French-language documents collected from
institutional, associative, and territorial sources, together with action sheets
provided by the SIAGE Chair. Its intended purpose is documentary exploration and
research support, rather than the profiling, ranking, or evaluation of persons
or organisations. Nevertheless, the construction of a searchable corpus and the
use of a language model raise data-governance, epistemic, and environmental
questions that must be considered throughout the pipeline.

\subsection{Data governance and reuse of documentary sources}

Most collected materials are publicly accessible documents describing projects,
public actions, or organisations. Public accessibility does not, however, remove
all legal and ethical obligations. A document may contain the name or contact
details of an individual, a testimony, or information that becomes identifying
when combined with other data. Under the GDPR, a personal datum is any
information relating to an identified or identifiable natural person; data
published online can therefore still be personal data \cite{cnil2024reuse}.

The project consequently follows a principle of data minimisation. It collects
documents because they describe initiatives relevant to territorial ageing, not
to construct personal profiles. The structured database retains the documentary
metadata needed for exploration---for example, source, territory, theme,
initiative holder, and date---but does not intentionally infer sensitive
attributes about individuals. Provenance is retained so that a user can return
to the original source, verify an extracted value, and request a correction or
removal where appropriate. Future public releases should also verify the terms
of reuse and copyright conditions for each source, respect website access
conditions, and avoid redistributing full texts when this is not permitted.

\subsection{Human oversight and responsible interpretation}

The model-generated records should be understood as navigational metadata, not
as authoritative factual statements. A language model may omit a value,
normalise an ambiguous expression too strongly, or produce an unsupported
interpretation. This risk is particularly important for fields such as the
initiative holder and date, for which source documents may not clearly separate
an organiser from a partner or a publication date from an initiative date.

Several design choices mitigate this risk. The extraction prompt instructs the
model not to invent information and to return an empty value when evidence is
insufficient. The original text, source identifier, prompt configuration, and,
where available, raw model response are retained to support auditability. The
human reference evaluation reported in Chapter~6 further distinguishes correct,
partial, incorrect, correctly absent, and non-evaluable cases. In any future
interface, uncertain metadata should remain linked to its source and should not
be presented as a verified claim without an appropriate review process.

The corpus itself also conditions what can be concluded from it. It is
French-language, source-selected, and strongly represented by VADA records.
Consequently, frequency analyses reflect both territorial initiatives and the
editorial choices of the sources. The controlled labels used for target audience
and territorial scale facilitate consistent indexing, but they simplify diverse
social situations. They must not be treated as exhaustive descriptions of
people's identities, needs, or experiences. The system should support expert
interpretation and exploratory access; it should not be used to make automated
decisions about funding, eligibility, priority, or the quality of an initiative.

\subsection{Local processing and environmental sobriety}

The project uses a local \texttt{qwen2.5:1.5b} model through Ollama for the
LLM-assisted extraction stage. Local processing limits the transfer of source
texts to an external commercial API and gives the project team control over the
inference environment. It does not, by itself, demonstrate a lower environmental
impact: this depends on the hardware, electricity mix, number of runs, and model
configuration. The project did not measure electricity consumption or carbon
emissions, so no quantitative footprint is claimed here.

Nevertheless, the implementation adopts several proportionate measures. It uses
a relatively small model, performs extraction on a defined corpus rather than
training a new model, uses deterministic decoding (temperature 0), caps generated
output at 700 tokens, and archives usable outputs to avoid unnecessary repeated
inference. Input texts are capped at 2,500 characters to bound computational
cost, although this introduces a documented recall risk when relevant information
appears late in a document. These choices are motivated by reproducibility and
operational constraints as well as by computational sobriety. They do not remove
the environmental cost of inference or of maintaining data infrastructure.

Environmental impact is a relevant NLP research concern because computational
choices entail energy and material costs \cite{strubell2019energy}. In line with
broader AI-ethics guidance, future versions should record the model version,
hardware, number of processed documents, number of inference runs, and, when
feasible, energy-use measurements \cite{unesco2021recommendation}. Such records
would make it possible to compare extraction quality with operational and
environmental cost rather than treating accuracy as the only optimisation target.

\subsection{Scope and future safeguards}

The current safeguards are appropriate to a research prototype, but they do not
replace governance for a public service. Before deployment beyond the internship,
the project should define a source-by-source reuse register, a correction and
removal procedure, retention rules for raw data and model outputs, and clear
responsibilities for validating public-facing metadata. A small second human
annotation would also help assess whether the reference labels are reproducible.
Together, these measures would preserve the main value of Index\_Insight---making
dispersed knowledge easier to find---while avoiding the presentation of automated
extraction as neutral, exhaustive, or error-free.

% ---------------------------------------------------------------------------
% BibTeX entries to add to your .bib file:
%
%@misc{cnil2024reuse,
%  author       = {{Commission nationale de l'informatique et des libert\'es}},
%  title        = {R\'eutilisation de vos donn\'ees publi\'ees sur Internet \`a
%                  des fins commerciales : quels sont vos droits ?},
%  year         = {2024},
%  month        = jun,
%  url          = {https://www.cnil.fr/fr/reutilisation-de-vos-donnees-publiees-sur-internet-des-fins-commerciales-quels-sont-vos-droits},
%  urldate      = {2026-08-12}
%}
%
%@inproceedings{strubell2019energy,
%  author       = {Strubell, Emma and Ganesh, Ananya and McCallum, Andrew},
%  title        = {Energy and Policy Considerations for Deep Learning in {NLP}},
%  booktitle    = {Proceedings of the 57th Annual Meeting of the Association for Computational Linguistics},
%  pages        = {3645--3650},
%  year         = {2019},
%  doi          = {10.18653/v1/P19-1355}
%}
%
%@misc{unesco2021recommendation,
%  author       = {{UNESCO}},
%  title        = {Recommendation on the Ethics of Artificial Intelligence},
%  year         = {2021},
%  month        = nov,
%  url          = {https://www.unesco.org/en/legal-affairs/recommendation-ethics-artificial-intelligence},
%  urldate      = {2026-08-12}
%}
% ---------------------------------------------------------------------------
